\subsection{Homomorphisms of a finitely presented module}

Let E be an A-module. If I is a filtered preordered set and \((\mathbf{G}_{i}, u_{ji})\) an inductive system of A-modules relative to I, the canonical maps \(G_{i} \to \varinjlim G_{i}\) induce homomorphisms \(\operatorname{Hom}_{\mathbf{A}}(\mathbf{E}, \mathbf{G}_{i}) \to \operatorname{Hom}_{\mathbf{A}}(\mathbf{E}, \varinjlim \mathbf{G}_{i})\) whence a homomorphism said to be canonical

\[
\varinjlim_ {i \in \mathrm{I}} \operatorname{Hom} _ {\mathrm{A}} (\mathrm{E}, \mathrm{G} _ {i}) \rightarrow \operatorname{Hom} _ {\mathrm{A}} (\mathrm{E}, \varinjlim_ {i \in \mathrm{I}} \mathrm{G} _ {i}).\tag{16}
\]

Let B be another ring, F a B-module, G an (A, B)-bimodule; we defined in II, p.~75 a canonical homomorphism:

\[
\operatorname{Hom} _ {\mathbf {A}} (\mathrm{E}, \mathrm{G}) \otimes_ {\mathbf {B}} \mathrm{F} \rightarrow \operatorname{Hom} _ {\mathbf {A}} (\mathrm{E}, \mathrm{G} \otimes_ {\mathbf {B}} \mathrm{F}).\tag{17}
\]

PROPOSITION 8. --- a) If the A-module E is finitely generated (resp. finitely presented), the canonical homomorphism (16) is injective (resp. bijective).

\begin{enumerate}
\def\labelenumi{\alph{enumi})}
\setcounter{enumi}{1}
\tightlist
\item
  Suppose that the B-module F is flat; if the A-module E is finitely generated (resp. finitely presented), the canonical homomorphism (17) is injective (resp. bijective). Let us prove, for example, b), the proof of a) being analogous. Consider A, B, F, G as fixed, and, for every right A-module E, set
\end{enumerate}

\[
\mathrm{T} (\mathrm{E}) = \operatorname{Hom} _ {\mathrm{A}} (\mathrm{E}, \mathrm{G}) \otimes_ {\mathrm{B}} \mathrm{F}, \quad \mathrm{T} ^ {\prime} (\mathrm{E}) = \operatorname{Hom} _ {\mathrm{A}} (\mathrm{E}, \mathrm{G} \otimes_ {\mathrm{B}} \mathrm{F})
\]

and let us denote by \(\nu_{\mathbf{E}}\) the homomorphism (17); for every homomorphism \(v: \mathbf{E} \to \mathbf{E}'\) of right A-modules, set \(\mathbf{T}(v) = \operatorname{Hom}(v, 1_{\mathbf{G}}) \otimes 1_{\mathbf{F}}\) and \(\mathbf{T}'(v) = \operatorname{Hom}(v, 1_{\mathbf{G}} \otimes 1_{\mathbf{F}})\) .

Let \(L_{1} \xrightarrow{v} L_{0} \xrightarrow{w} E \to 0\) be a presentation of E; we assume the free module \(L_{0}\) (resp. the free modules \(L_{0}\) and \(L_{1}\) ) to be finitely generated. The diagram

\[
\begin{array}{c} 0 \to \mathrm{T(E)} \xrightarrow {\mathrm{T(w)}} \mathrm {T(L_ {0})} \xrightarrow {\mathrm{T(v)}} \mathrm {T(L_ {1})} \\ v _ {\mathrm{E}} \Big \downarrow \qquad \qquad \qquad \qquad \qquad \qquad \qquad \qquad \qquad \qquad \qquad \qquad \qquad \qquad \qquad \qquad \qquad \qquad \qquad \qquad \qquad \qquad \qquad \qquad \qquad \qquad \qquad \qquad \qquad \qquad \qquad \qquad \qquad \qquad 0 \to \mathrm {T^ {\prime} (E)} \xrightarrow [ \mathrm {T^ {\prime} (w)} ]{} \mathrm {T^ {\prime} (L_ {0})} \xrightarrow [ \mathrm {T^ {\prime} (v)} ]{} \mathrm {T^ {\prime} (L_ {1})} \end{array}\tag{18}
\]

is commutative, and its second row is exact (II, p.~36, th. 1); moreover, the sequence

\[
0 \rightarrow \operatorname{Hom} _ {\mathbf {A}} (\mathrm{E}, \mathrm{G}) \rightarrow \operatorname{Hom} _ {\mathbf {A}} (\mathrm{L} _ {0}, \mathrm{G}) \rightarrow \operatorname{Hom} _ {\mathbf {A}} (\mathrm{L} _ {1}, \mathrm{G})
\]

is exact (loc. cit.), and since F is flat, the first row of (18) is also an exact sequence (X, p.~8, def. 1). That being so, we know that \(v_{L_0}\) is bijective (resp. that \(v_{L_0}\) and \(v_{L_1}\) are bijective) (II, p.~75, prop. 2). If one assumes only \(v_{L_0}\) bijective, it follows from (18) that \(v_{L_0} \circ T(w) = T'(w) \circ v_E\) is injective, therefore \(v_E\) is also. If we suppose that \(v_{L_0}\) and \(v_{L_1}\) both are bijective, we deduce from Cor. 2 (ii) of X, p.~6 that \(v_E\) is surjective, and since we have just seen that \(v_E\) is injective, it is bijective.

COROLLARY. --- Every flat module of finite presentation is projective.

Indeed, let E be a flat A-module of finite presentation. Applying (b) to the case \(\mathbf{B} = \mathbf{A}\) , \(\mathbf{G} = {}_{\mathrm{s}}\mathbf{A}_{\mathrm{d}}\) , \(\mathbf{F} = \mathbf{E}\) , we see that the canonical homomorphism

\[
\operatorname{Hom} _ {\mathbf {A}} (\mathbf {E}, \mathbf {A}) \otimes_ {\mathbf {A}} \mathbf {E} \rightarrow \operatorname{Hom} _ {\mathbf {A}} (\mathbf {E}, \mathbf {E})
\]

is surjective. This implies that E is projective (II, p.~77, Remark 1).

By the preceding corollary and Prop. 5 of X, p.~10, there is an identity between finitely presented flat modules and finitely generated projective modules. On the other hand, there exist finitely generated flat modules that are not finitely presented, hence not projective (cf.~X, p.~170, Exercise 17; see however X, p.~169, Exercises 13 and 14).

